\RequirePackage[l2tabu, orthodox]{nag}
\RequirePackage{silence}
\WarningFilter{fmtcount}{\ordinal already defined use \FCordinal instead}
\documentclass[french, english]{beamer}
%INSTALL

%avoids a warning
\usepackage[log-declarations=false]{xparse}
\usepackage{fontspec} %font selecting commands
\usepackage{xunicode}
%warn about missing characters
\tracinglostchars=2

%REDAC
\usepackage{booktabs}
\usepackage{calc}

\usepackage{mathtools} %load this before babel!
	\mathtoolsset{showonlyrefs,showmanualtags}

\usepackage{babel}
%suppresses the warning about frenchb not modifying the captions (“—” to “:” in “Figure 1 – Legend”).
	\frenchbsetup{AutoSpacePunctuation=false,SuppressWarning=false}

\usepackage[super]{nth}
\usepackage{listings} %typeset source code listings
	\lstset{language=XML,tabsize=2,literate={"}{{\tt"}}1,captionpos=b}
\usepackage[nolist,smaller,printonlyused]{acronym}%,smaller option produces warnings from relsize in some cases, it seems.
\usepackage[nodayofweek]{datetime}%must be loaded after the babel package
\usepackage{xspace}
\usepackage{hyperref}% option pdfusetitle must be introduced here, not in hypersetup.
%breaklinks makes links on multiple lines into different PDF links to the same target.
%colorlinks (false): Colors the text of links and anchors. The colors chosen depend on the the type of link. In spite of colored boxes, the colored text remains when printing.
%linkcolor=black: this leaves other links in colors, e.g. refs in green, don't print well.
%pdfborder (0 0 1, set to 0 0 0 if colorlinks): width of PDF link border
%hidelinks
\hypersetup{breaklinks,bookmarksopen,colorlinks=true,urlcolor=blue,linkcolor=,hyperfigures=true}
% hyperref doc says: Package bookmark replaces hyperref’s bookmark organization by a new algorithm (...) Therefore I recommend using this package.
\usepackage{bookmark}

% center floats by default, but do not use with float
% \usepackage{floatrow}
% \makeatletter
% \g@addto@macro\@floatboxreset\centering
% \makeatother
\usepackage{ragged2e} %new com­mands \Cen­ter­ing, \RaggedLeft, and \RaggedRight and new en­vi­ron­ments Cen­ter, FlushLeft, and FlushRight, which set ragged text and are eas­ily con­fig­urable to al­low hy­phen­ation (the cor­re­spond­ing com­mands in LaTeX, all of whose names are lower-case, pre­vent hy­phen­ation al­to­gether). 
\usepackage{siunitx} %[expproduct=tighttimes, decimalsymbol=comma]
\sisetup{detect-all}% to detect e.g. when in math mode (use a math font)
\usepackage{braket} %for \Set
\usepackage{natbib}

\usepackage{amsmath,amsthm}
% \usepackage{amsfonts} %not required?
% \usepackage{dsfont} %for what?
%unicode-math overwrites the following commands from the mathtools package: \dblcolon, \coloneqq, \Coloneqq, \eqqcolon. Using the other colon-like commands from mathtools will lead to inconsistencies. Plus, Using \overbracket and \underbracke from mathtools package. Use \Uoverbracket and \Uunderbracke for original unicode-math definition.
%use exclusively \mathbf and choose math bold style below.
\usepackage[warnings-off={mathtools-colon, mathtools-overbracket}, bold-style=ISO]{unicode-math}

\defaultfontfeatures{
	Fractions=On,
	Mapping=tex-text% to turn "--" into dashes, useful for bibtex%%
}
\defaultfontfeatures[\rmfamily, \sffamily]{
	Fractions=On,
	Mapping=% to leave " alone (disable the default mapping tex-text; requires loading the font afterwards?)
}
\newfontfamily\xitsfamily{XITS}
\newfontfamily\texgyretermesfamily{TeX Gyre Termes}
\newfontfamily\texgyreherosfamily{TeX Gyre Heros}
\newfontfamily\lmfamily{Latin Modern Roman}
\setmainfont{Latin Modern Roman}
\setsansfont{Latin Modern Sans}
\setmonofont{Latin Modern Mono}
\defaultfontfeatures{}%disable default font features to avoid warnings with math fonts.
\setmathfont{XITS Math}
\setmathfont[range={\mathcal,\mathbfcal},StylisticSet=1]{XITS Math}

\usepackage{cleveref}% cleveref should go "laster" than hyperref
%GRAPHICS
\usepackage{pgf}
\usepackage{pgfplots}
	\usetikzlibrary{matrix,fit,plotmarks,calc,trees,shapes.geometric,positioning,plothandlers}
\pgfplotsset{compat=1.11}
\usepackage{graphicx}

\graphicspath{{graphics/},{graphics-dm/}}
\DeclareGraphicsExtensions{.pdf}

%HACKING
\usepackage{printlen}
\uselengthunit{mm}
% 	\newlength{\templ}% or LenTemp?
% 	\setlength{\templ}{6 pt}
% 	\printlength{\templ}
\usepackage{etoolbox} %for addtocmd command
\usepackage{scrhack}% load at end. Corrects a bug in float package, which is outdated but might be used by other packages
\usepackage{xltxtra} %somebody said that this is loaded by fontspec, but does not seem correct: if not loaded explicitly, does not appear in the log and \showhyphens is not corrected.

%Beamer-specific
%ADD
\usepackage{appendixnumberbeamer}
%\setbeamersize{text margin left=0.1cm, text margin right=0.1cm} 
\setbeamertemplate{navigation symbols}{}
\usetheme{BrusselsBelgium}
\usefonttheme{professionalfonts}


\newcommand{\R}{ℝ}
\newcommand{\N}{ℕ}
\newcommand{\Z}{ℤ}
\newcommand{\card}[1]{\lvert{#1}\rvert}
\newcommand{\powerset}[1]{\mathscr{P}(#1)}%\mathscr rather than \mathcal: scr is rounder than cal (at least in XITS Math).
\newcommand{\suchthat}{\;\ifnum\currentgrouptype=16 \middle\fi|\;}
%\newcommand{\Rplus}{\reels^+\xspace}

\AtBeginDocument{%
	\renewcommand{\epsilon}{\varepsilon}
% we want straight form of \phi for mathematics, as recommended in UTR #25: Unicode support for mathematics.
%	\renewcommand{\phi}{\varphi}
}

% with amssymb, but I don’t want to use amssymb just for that.
% \newcommand{\restr}[2]{{#1}_{\restriction #2}}
%\newcommand{\restr}[2]{{#1\upharpoonright}_{#2}}
\newcommand{\restr}[2]{{#1|}_{#2}}%sometimes typed out incorrectly within \set.
%\newcommand{\restr}[2]{{#1}_{\vert #2}}%\vert errors when used within \Set and is typed out incorrectly within \set.
\DeclareMathOperator*{\argmax}{arg\,max}
\DeclareMathOperator*{\argmin}{arg\,min}


%ARG TH
\newcommand{\AF}{\mathcal{AF}}
\newcommand{\labelling}{\mathcal{L}}
\newcommand{\labin}{\textbf{in}\xspace}
\newcommand{\labout}{\textbf{out}}
\newcommand{\labund}{\textbf{undec}\xspace}
\newcommand{\nonemptyor}[2]{\ifthenelse{\equal{#1}{}}{#2}{#1}}
\newcommand{\gextlab}[2][]{
	\labelling{\mathcal{GE}}_{(#2, \nonemptyor{#1}{\ibeatsr{#2}})}
}
\newcommand{\allargs}{A^*}
\newcommand{\args}{A}
\newcommand{\ar}{a}

%MCDA+Arg
\newcommand{\dm}{d}
\newcommand{\ileadsto}{\leadsto}
\newcommand{\mleadsto}[1][\eta]{\leadsto_{#1}}
\newcommand{\ibeats}{\vartriangleright}
\newcommand{\mbeats}[1][\eta]{\vartriangleright_{#1}}


%MISC
\newcommand{\lequiv}{\Vvdash}
\newcommand{\weightst}{W^{\,t}}

%MCDA classical
\newcommand{\crits}{\mathcal{J}}
\newcommand{\altspace}{\mathbb{A}}
\newcommand{\alts}{A}

%Sorting
\newcommand{\cats}{\mathcal{C}}
\newcommand{\catssubsets}{2^\cats}
\newcommand{\catgg}{\vartriangleright}
\newcommand{\catll}{\vartriangleleft}
\newcommand{\catleq}{\trianglelefteq}
\newcommand{\catgeq}{\trianglerighteq}
\newcommand{\alttoc}[2][x]{(#1 \xrightarrow{} #2)}
\newcommand{\alttocat}[3]{(#2 \xrightarrow{#1} #3)}
\newcommand{\alttoI}{(x \xrightarrow{} \left[\underline{C_x}, \overline{C_x}\right])}
\newcommand{\alttocatdm}[3][t]{\left(#2 \thinspace \raisebox{-3pt}{$\xrightarrow{#1}$}\thinspace #3\right)}
\newcommand{\alttocatatleast}[2]{\left(#1 \thinspace \raisebox{-3pt}{$\xrightarrow[]{≥}$}\thinspace #2\right)}
\newcommand{\alttocatatmost}[2]{\left(#1 \thinspace \raisebox{-3pt}{$\xrightarrow[]{≤}$}\thinspace #2\right)}

\newcommand{\source}{\scriptsize}
\newcommand{\commentOC}[1]{{\selectlanguage{french}{\todo{OC : #1}}}}
%Or: \todo[color=green!40]

%this probably requires outdated float package, see doc KomaScript for an alternative.
% \newfloat{program}{t}{lop}
% \floatname{program}{PM}

%\crefname{axiom}{axiom}{axioms}%might be needed for workaround bug in cref when defining new theorems?

%\ifdefined\theorem\else
%\newtheorem{theorem}{\iflanguage{english}{Theorem}{Théorème}}
%\fi

%which line breaks are chosen: accept worse lines, therefore reducing risk of overfull lines. Default = 200
\tolerance=2000
%accept overfull hbox up to...
\hfuzz=2cm
%reduces verbosity about the bad line breaks
\hbadness 5000
%sloppy sets tolerance to 9999
\apptocmd{\sloppy}{\hbadness 10000\relax}{}{}

% WRITING
%\newcommand{\ie}{i.e.\@\xspace}%to try
%\newcommand{\eg}{e.g.\@\xspace}
%\newcommand{\etal}{et al.\@\xspace}
\newcommand{\ie}{i.e.\ }
\newcommand{\eg}{e.g.\ }
\newcommand{\mkkOK}{\checkmark}%\color{green}{\checkmark}
\newcommand{\mkkREQ}{\ding{53}}%requires pifont?%\color{green}{\checkmark}
\newcommand{\mkkNO}{}%\text{\color{red}{\textsf{X}}}

\makeatletter
\newcommand{\boldor}[2]{%
	\ifnum\strcmp{\f@series}{bx}=\z@
		#1%
	\else
		#2%
	\fi
}
\newcommand{\textstyleElProm}[1]{\boldor{\MakeUppercase{#1}}{\textsc{#1}}}
\makeatother
\newcommand{\electre}{\textstyleElProm{Électre}\xspace}
\newcommand{\electreIv}{\textstyleElProm{Électre Iv}\xspace}
\newcommand{\electreIV}{\textstyleElProm{Électre IV}\xspace}
\newcommand{\electreIII}{\textstyleElProm{Électre III}\xspace}
\newcommand{\electreTRI}{\textstyleElProm{Électre Tri}\xspace}
% \newcommand{\utadis}{\texorpdfstring{\textstyleElProm{utadis}\xspace}{UTADIS}}
% \newcommand{\utadisI}{\texorpdfstring{\textstyleElProm{utadis i}\xspace}{UTADIS I}}

%TODO
% \newcommand{\textstyleElProm}[1]{{\rmfamily\textsc{#1}}} 


\newlength{\GraphsNodeSep}
\setlength{\GraphsNodeSep}{7mm}

% MCDA Drawing Sorting
\newlength{\MCDSCatHeight}
\setlength{\MCDSCatHeight}{6mm}
\newlength{\MCDSAltHeight}
\setlength{\MCDSAltHeight}{4mm}
%separation between two vertical alts
\newlength{\MCDSAltSep}
\setlength{\MCDSAltSep}{2mm}
\newlength{\MCDSCatWidth}
\setlength{\MCDSCatWidth}{3cm}
\newlength{\MCDSEvalRowHeight}
\setlength{\MCDSEvalRowHeight}{6mm}
\newlength{\MCDSAltsToCatsSep}
\setlength{\MCDSAltsToCatsSep}{1.5cm}
\newcounter{MCDSNbAlts}
\newcounter{MCDSNbCats}
\newlength{\MCDSArrowDownOffset}
\setlength{\MCDSArrowDownOffset}{0mm}

\tikzset{/Graphs/dot/.style={
	shape=circle, fill=black, inner sep=0, minimum size=1mm
}}
\tikzset{/MC/D/S/alt/.style={
	shape=rectangle, draw=black, inner sep=0, minimum height=\MCDSAltHeight, minimum width=2.5cm, anchor=north east
}}
\tikzset{MC/D/S/pref/.style={
	shape=ellipse, draw=gray, thick
}}
\tikzset{/MC/D/S/cat/.style={
	shape=rectangle, draw=black, inner sep=0, minimum height=\MCDSCatHeight, minimum width=\MCDSCatWidth, anchor=north west
}}
\tikzset{/MC/D/S/evals matrix/.style={
	matrix, row sep=-\pgflinewidth, column sep=-\pgflinewidth, nodes={shape=rectangle, draw=black, inner sep=0mm, text depth=0.5ex, text height=1em, minimum height=\MCDSEvalRowHeight, minimum width=12mm}, nodes in empty cells, matrix of nodes, inner sep=0mm, outer sep=0mm, row 1/.style={nodes={draw=none, minimum height=0em, text height=, inner ysep=1mm}}
}}

% Beliefs
\tikzset{/Beliefs/D/S/attacker/.style={
	shape=rectangle, draw, minimum size=8mm
}}
\tikzset{/Beliefs/D/S/supporter/.style={
	shape=circle, draw
}}

\newcommand{\tikzmark}[1]{%
	\tikz[overlay, remember picture, baseline=(#1.base)] \node (#1) {};%
}


\begin{acronym}
\acro{AMCD}{Aide Multicritère à la Décision}
\acro{ASA}{Argument Strength Assessment}
\acro{DA}{Decision Analysis}
\acro{DM}{Decision Maker}
\acro{DPr}{Deliberated Preferences}
\acro{DRSA}{Dominance-based Rough Set Approach}
\acro{DSS}{Decision Support Systems}
\acrodefplural{DSS}{Decision Support Systems}
% \newacroplural{DSS}[DSSes]{Decision Support Systems}
\acro{EJOR}{European Journal of Operational Research}
\acro{LNCS}{Lecture Notes in Computer Science}
\acro{MCDA}{Multicriteria Decision Aid}
\acro{MIP}{Mixed Integer Program}
\acro{NCSM}{Non Compensatory Sorting Model}
\acro{PL}{Programme Linéaire}
\acro{PLNE}{Programme Linéaire en Nombres Entiers}
\acro{PM}{Programme Mathématique}
\acro{MP}{Mathematical Program}
\acro{MIP}{Mixed Integer Program}
% \newacroplural{PM}{Programmes Mathématiques}
%acrodefplural since version 1.35, my debian has \ProvidesPackage{acronym}[2009/01/25, v1.34, Support for acronyms (Tobias Oetiker)]
\acrodefplural{PM}{Programmes Mathématiques}
\acro{PMML}{Predictive Model Markup Language}
\acro{RESS}{Reliability Engineering \& System Safety}
\acro{SMAA}{Stochastic Multicriteria Acceptability Analysis}
\acro{URPDM}{Uncertainty and Robustness in Planning and Decision Making}
\acro{XML}{Extensible Markup Language}
\end{acronym}

%this approach does not generalize to multipart nodes.
%\tikzset{/uml/interface/.style={rectangle, draw, fill=yellow!20, align=left, node font=\fontspec{Latin Modern Mono Light}, node contents={$<<$interface$>>$\\\textbf{#1}}, name=#1
%}}
%this approach does not apply for <<interface>>. I’ll do it manually, as I can’t find anything better. Use: \nodepart[font=]{one} \small <<interface>>\\\bfseries ItemDAO
\tikzset{/uml/class/.style={rectangle, draw, align=center, fill=yellow!20, node font=\ttfamily, font=\bfseries
}}
\tikzset{/uml/abstract class/.style={rectangle, draw, align=center, fill=yellow!20, node font=\fontspec{Latin Modern Mono Light}, font=\bfseries\itshape
}}
\tikzset{/uml/class3/.style={
rectangle split, rectangle split every empty part={}, rectangle split parts=3, rectangle split part align={center, left, left}, draw, fill=yellow!20, rectangle split empty part height=0, node font=\ttfamily, every one node part/.style={font=\bfseries, align=center}, every two node part/.style={align=left}, every three node part/.style={align=left}
}}
\tikzset{/uml/extends/.style={draw, -open triangle 45}}
\tikzset{/uml/implements/.style={draw, -open triangle 45, dashed}}
%prefix= used to cancel \texttt, otherwise this cancels the effect of the font selection \fontspec{Latin Modern Mono Light}, and thus renders \bfseries inoperant. With prefix=, the result is correct.
%\jeeref[prefix=]{javax.faces.component.html/HtmlCommandButton}\\
\tikzset{/uml/table/.style={rectangle, draw, align=center, fill=blue!20, node font=\ttfamily, font=\bfseries
}}
\tikzset{/uml/table2/.style={
rectangle split, rectangle split every empty part={}, rectangle split parts=2, rectangle split part align={center, left}, draw, fill=blue!20, rectangle split empty part height=0, node font=\ttfamily, every one node part/.style={font=\bfseries, align=center}, every two node part/.style={align=left}
}}
\tikzset{/uml/dbkey/.style={draw, ->}}
%http://tex.stackexchange.com/questions/79781/placing-anchor-before-and-after-text-in-multipart-rectangle



\newbool{refAPIIsAnnotation}
\newbool{refAPIDisplayFull}
\newbool{refAPIOmitFQName}
\pgfkeys{
	/refAPI/.is family,
	/refAPI/.cd,
	base url/.estore in = \refAPIBaseUrl,
	@/.is if = refAPIIsAnnotation,
	full/.is if = refAPIDisplayFull,
	omitfqname/.is if = refAPIOmitFQName,
	prefix/.store in = \refAPIPrefix,
	suffix/.estore in = \refAPISuffix,
	default/.style = {
		base url = ,
		@ = false,
		full = false,
		omitfqname = false,
		prefix = \texttt,
		suffix = ,
	},
}

%Usage: \refAPIParseRef{javax.persistence/EntityManager} ; \refAPIParseRef{javax.persistence/PersistenceContextType\#EXTENDED}
%Defines: \refAPIFQName (before #) ; \refAPIFieldFull (after #) ; \refAPIFieldShort (after # and before -) ; \refAPIPackageSlashes (FQName before / with slashes instead of dots) ; \refAPIClass (FQName after /)
\newcommand{\refAPIParseRef}[1]{%
	\IfSubStr{#1}{\#}{%
		\StrBefore{#1}{\#}[\refAPIFQName]%
		\StrBehind{#1}{\#}[\refAPIFieldFull]%
	}{%
		\edef\refAPIFQName{#1}%
		\edef\refAPIFieldFull{}%
	}%
	\IfSubStr{\refAPIFieldFull}{-}{%
		\StrBefore{\refAPIFieldFull}{-}[\refAPIFieldShort]%
	}{%
		\edef\refAPIFieldShort{\refAPIFieldFull}%
	}%
	\StrBefore{\refAPIFQName}{/}[\refAPIPackage]%
	\StrBehind{\refAPIFQName}{/}[\refAPIClass]%
	\StrSubstitute{\refAPIPackage}{.}{/}[\refAPIPackageSlashes]%
}

%Usage: \jeeref{javax.persistence/EntityManager} ; \jeeref[@]{javax.persistence/PersistenceContextType\#EXTENDED}
\newcommand{\jeeref}[2][]{\refAPIParseRef{#2}\jrefAPI{base url = https://docs.oracle.com/javaee/7/api/, #1}}
\newcommand{\jseref}[2][]{\refAPIParseRef{#2}\jrefAPI{base url = https://docs.oracle.com/javase/8/docs/api/, #1}}
\newcommand{\jsftag}[2][]{\refAPIParseRef{#2}\jrefAPI{base url = https://docs.oracle.com/javaee/7/javaserver-faces-2-2/vdldocs-facelets/, full, prefix=:, #1}}
%Standard HTML RenderKit : https://docs.oracle.com/javaee/7/javaserver-faces-2-2/renderkitdocs/HTML_BASIC/javax.faces.Inputjavax.faces.Text.html

%Accepts keys: @, prefix, suffix. Requires key: base url. Expects defined: \refAPIField ; \refAPIPackageSlashes ; \refAPIClass
%Usage: \texttt{"Hello".\jseref[omitfqname]{java.lang/String\#substring-int-int-}(1, 3);}
\newcommand{\jrefAPI}[1]{%
	\pgfkeys{/refAPI/.cd, default, #1}%
	\ifbool{refAPIIsAnnotation}{%
		\edef\refAPIAnnot{@}%
	}{%
		\edef\refAPIAnnot{}%
	}%
	\IfEq{\refAPIFieldFull}{}{%
		{%
			\ifbool{refAPIDisplayFull}{%
				\refAPIPackage%
			}{%
			}%
			\refAPIPrefix{%
			\href{%
				\refAPIBaseUrl\refAPIPackageSlashes/\refAPIClass.html%
			}{%
				\refAPIAnnot\refAPIClass%
			}%
		}\refAPISuffix}%
	}{%
		{\refAPIPrefix{%
			\ifbool{refAPIOmitFQName}{%
			}{%
				\refAPIAnnot\refAPIClass.%
			}%
			\href{%
				\refAPIBaseUrl\refAPIPackageSlashes/\refAPIClass.html\#\refAPIFieldFull%
			}{%
				\refAPIFieldShort%
			}%
		}\refAPISuffix}%
	}%
}



\setbeamertemplate{headline}[singleline]
\setbeamertemplate{footline}[onlypage]

\title{SWT}
\subtitle{}
\subject{Java}
\keywords{GUI}
\author{Olivier Cailloux}
\institute[LAMSADE]{LAMSADE, Université Paris-Dauphine}
\date{Version of \today}

\begin{document}
\bibliographystyle{apalike}

\begin{frame}[plain]
	\tikz[remember picture,overlay]{
		\path (current page.south west) node[anchor=south west, inner sep=0] {
			\includegraphics[height=1cm]{LAMSADE95.jpg}
		};
		\path (current page.south) ++ (0, 1mm) node[anchor=south, inner sep=0] {
			\includegraphics[height=9mm]{Dauphine.jpg}
		};
		\path (current page.south east) node[anchor=south east, inner sep=0] {
			\includegraphics[height=1cm]{PSL.png}
		};
	}
   \titlepage
\end{frame}
\addtocounter{framenumber}{-1}

\section{GUIs in general}
\begin{frame}
	\frametitle{What’s a GUI?}
	\begin{itemize}
		\item GUI? \pause Graphical User Interface
		\item \emph{An} interface between the user and your program
		\item Assembles graphical components
		\item Components are provided
		\item By the OS or language-specific
		\item You can design your own components
		\item Own components: rare!
	\end{itemize}
\end{frame}

\begin{frame}
	\frametitle{Control in GUIs}
	\begin{itemize}
		\item IoC? \pause Inversion of control
		\item AKA The Hollywood principle: “Don’t call us, we’ll call you”
		\item Usually: you control the application flow
		\item From \texttt{main}, go to this method, …
		\item With a GUI, the GUI controls the flow
		\item Most of the time: do “nothing”
		\begin{itemize}
			\item Meaning: nothing you program
			\item Refresh components when necessary, …
			\item Done by the GUI library
		\end{itemize}
		\item Your program acts when the user clicks a button, …
		\item Called an event
	\end{itemize}
\end{frame}

\begin{frame}[fragile]
	\frametitle{Hierarchical design}
	Components are {\tiny typically} assembled into a tree
	\begin{example}[A tree of GUI components]
		\begin{tikzpicture}
			\path node[matrix of nodes, row sep=5mm] {
				& & & & |(App)| App\\
				& & & |(Win1)| Win1 & & |(Win2)| Win2\\
				& |(Left)| Left pane & |(Main)| Main pane & |(Menu)| Menu bar & |(Status)| Status bar\\
				|(Top)| Top text & & |(Middle)| Left content\\
			};
				
			\path graph[use existing nodes, left anchor=south, right anchor=north] {
				App -- {
					Win1 -- {
						Left -- {Top, Middle},
						Main,
						Menu,
						Status
					}, 
					Win2
				}
			}; 
		\end{tikzpicture}
	\end{example}
\end{frame}

\begin{frame}
	\frametitle{Observer situation}
	Situation: One object must advertise events to other objects
	\begin{example}[Advertiser]
		\begin{itemize}
			\item A button advertises that it has been pressed
			\item A data store warns that it has almost no space left
		\end{itemize}
	\end{example}
	Simplest approach: 
	\begin{itemize}
		\item The advertiser knows the interested objects and calls them
	\end{itemize}
	\begin{center}
		\begin{tikzpicture}
			\path node[/uml/class3] (a) {
				\nodepart{one}
				AlmostFullHandler
				\nodepart{three}
				almostFullEvent()
			};
			\path (a.east) ++ (1cm, 0) node[anchor=west, /uml/class3] (d) {
				\nodepart{one}
				DataStore
				\nodepart{two}
				almostFullH
				\nodepart{three}
				warnAlmostFull()
			};
			\path (d.three south) ++ (0, -4mm) node[anchor=north] {\texttt{\{almostFullH.almostFullEvent();\}}} edge[-] (d.three east);
		\end{tikzpicture}
	\end{center}
\end{frame}

\begin{frame}
	\frametitle{Observer pattern}
	Situation:
	\begin{itemize}
		\item You want one object to advertise events to observers
		\item You do not know all possible observers at development time
		\item You want loose coupling between objects
	\end{itemize}
	\begin{itemize}
		\item The observed object calls any observer object
		\item Observer is an interface
		\item Observers must register to the observed object
	\end{itemize}
	\begin{center}
		\begin{tikzpicture}
			\path node[/uml/class3] (o) {
				\nodepart[font=]{one} \small <<interface>>\\\bfseries
				Observer
				\nodepart{three}
				almostFullEvent()
			};
			\path (o.east) ++ (1cm, 0) node[anchor=west, /uml/class3] (d) {
				\nodepart{one}
				DataStore
				\nodepart{two}
				observer
				\nodepart{three}
				registerObserver(o: Observer)\\
				warnAlmostFull()
			};
			\path (d.three south) node[anchor=north] {\texttt{\{o.almostFullEvent();\}}};
		\end{tikzpicture}
	\end{center}
\end{frame}

\begin{frame}
	\frametitle{Programming a GUI}
	Simple approach:
	\begin{itemize}
		\item Assemble your graphical components into a tree
		\item Register observers on components
		\item Start the GUI event loop
	\end{itemize}
\end{frame}

\begin{frame}
	\frametitle{GUIs in Java}
	\begin{itemize}
		\item Abstract Window Toolkit (AWT): oldest and most basic
		\item Uses native GUI components
		\item Swing: higher-level controls
		\item Reuses part of AWT to plug into OS
		\item Paints its own controls (consistent look and feel across OSes)
		\item JavaFX: replacer of Swing
		\item Declarative language for definition of the GUI; animations; CSS styles
		\item Paints its own controls
		\item AWT, Swing, JavaFX part of the JRE
		\item Standard Widget Toolkit (SWT): Eclipse’s alternative to AWT/Swing
		\item JFace: Eclipse’s advanced library that supplements SWT
	\end{itemize}
\end{frame}

\section{SWT}
\begin{frame}
	\frametitle{Characteristics}
	\begin{itemize}
		\item A blatant lie in the name (? \pause It’s not standard!)
		\item JavaFX is newer and (possibly) more popular, and (really) standard
		\item SWT uses native graphical components (like AWT)
		\item JFace/SWT has high-level components (like Swing or JavaFX)
		\item[⇒] Native look and feel with high-level components
	\end{itemize}
\end{frame}

\begin{frame}[fragile]
	\frametitle{Hierarchical structure}
	Example:
	\footnotesize
	\begin{tikzpicture}
		\path node[matrix of nodes, row sep=5mm] {
			& & & & |(App)| display: \texttt{Display}\\
			& & & |(Win1)| win1: \texttt{Shell} & & |(Win2)| win2\\
			& |(Left)| Left pane & |(Main)| Main pane: \texttt{Composite} & |(Menu)| Menu bar & |(Status)| Status bar\\
			|(Top)| Top text & & |(Middle)| Left content: \texttt{Tree} \\
		};
			
		\path graph[use existing nodes, left anchor=south, right anchor=north] {
			App -- {
				Win1 -- {
					Left -- {Top, Middle},
					Main,
					Menu,
					Status
				}, 
				Win2
			}
		}; 
	\end{tikzpicture}
\end{frame}

\begin{frame}
	\frametitle{Core concepts}
	\begin{minipage}[t]{\widthof{“Child of” relation}+3mm}
		\begin{block}{“Child of” relation}
			\begin{center}
				\begin{tikzpicture}
					\path node[/uml/class] (widget) {Widget};
					\path (widget.south) ++(0, -3mm) node[/uml/class, anchor=north] (control) {Control};
					\path (control.south) ++(0, -3mm) node[/uml/class, anchor=north] (composite) {Composite};
					\path (composite.south) ++(0, -3mm) node[/uml/class, anchor=north] (shell) {Shell};
					\path[/uml/extends] (shell) -- (composite);
					\path[/uml/extends] (composite) -- (control);
					\path[/uml/extends] (control) -- (widget);
				\end{tikzpicture}
			\end{center}
		\end{block}
	\end{minipage}\hfill
	\begin{minipage}[t]{\textwidth-35mm}
		\begin{itemize}
			\item \href{http://help.eclipse.org/neon/advanced/content.jsp?topic=/org.eclipse.platform.doc.isv/reference/api/org/eclipse/swt/widgets/Widget.html}{\texttt{Widget}}: GUI component; can send events; must be disposed
			\item \href{http://help.eclipse.org/neon/advanced/content.jsp?topic=/org.eclipse.platform.doc.isv/reference/api/org/eclipse/swt/widgets/Control.html}{\texttt{Control}}: may be put into a composite; has bounds
			\item \href{http://help.eclipse.org/neon/advanced/content.jsp?topic=/org.eclipse.platform.doc.isv/reference/api/org/eclipse/swt/widgets/Composite.html}{\texttt{Composite}}: can contain child controls
			\item \href{http://help.eclipse.org/neon/advanced/content.jsp?topic=/org.eclipse.platform.doc.isv/reference/api/org/eclipse/swt/widgets/Shell.html}{\texttt{Shell}}: top-level window connected to a display
			\item \href{http://help.eclipse.org/neon/advanced/content.jsp?topic=/org.eclipse.platform.doc.isv/reference/api/org/eclipse/swt/widgets/Display.html}{\texttt{Display}}: represents the connection to the native display
		\end{itemize}
	\end{minipage}
\end{frame}

\begin{frame}[fragile]
	\frametitle{General structure}
	\begin{itemize}
		\item Create the display
		\item Create a shell bound to the display
		\item Put the desired controls into the shell
		\item Run the event loop
	\end{itemize}
	\begin{lstlisting}
Display display = new Display();
Shell shell = new Shell(display);
// … populate shell …
shell.open();
while (!shell.isDisposed()) {
  if (!display.readAndDispatch()) {
    display.sleep();
  }
}
display.dispose();
	\end{lstlisting}
\end{frame}

\begin{frame}
	\frametitle{Create and position a control}
	\begin{itemize}
		\item Use its constructor
		\item Provide the parent composite, and the style bits
		\item \texttt{new Button(shell, SWT.PUSH | SWT.WRAP);}
		\item A control has a position (relative to its parent) and a size
		\item Bounds of the control: rectangle that contains it entirely (> client area)
		\item \texttt{control.setBounds(10, 0, 50, 20);} (location: (10, 0); size: 50 × 20)
		\item Hard to determine the right size! Depends on text content, …
		\item Compute preferred size: \texttt{control.computeSize(…)}
		\item Set to preferred size: \texttt{control.pack();}
	\end{itemize}
\end{frame}

\begin{frame}
	\frametitle{Listen to widgets}
	\begin{itemize}
		\item Implement the right *Listener interface
		\item Example: \href{http://help.eclipse.org/neon/advanced/content.jsp?topic=/org.eclipse.platform.doc.isv/reference/api/org/eclipse/swt/events/SelectionListener.html}{\texttt{SelectionListener}}
		\item Add your implementer to the observers of the widget
		\item \texttt{button.addSelectionListener(myListener);}
		\item When the widget fires the corresponding event, your listener will be called
	\end{itemize}
\end{frame}

\begin{frame}
	\frametitle{Dispose resources}
	\begin{itemize}
		\item Widgets relate to native resources
		\item Must be explicitly disposed of
		\item Rule: if you create it, you must dispose of it
		\item \texttt{composite.dispose()} ⇒ disposes of all the children
	\end{itemize}
\end{frame}

\begin{frame}
	\frametitle{Items in controls}
	\begin{itemize}
		\item Some controls must be populated with items
		\item Example: \href{http://help.eclipse.org/neon/advanced/content.jsp?topic=/org.eclipse.platform.doc.isv/reference/api/org/eclipse/swt/widgets/TableItem.html}{\texttt{TableItem}} represents a row of a \href{http://help.eclipse.org/neon/advanced/content.jsp?topic=/org.eclipse.platform.doc.isv/reference/api/org/eclipse/swt/widgets/Table.html}{\texttt{Table}}
		\item \href{http://help.eclipse.org/neon/advanced/content.jsp?topic=/org.eclipse.platform.doc.isv/reference/api/org/eclipse/swt/widgets/Item.html}{\texttt{Item}}s are widgets but not controls: their position is {\tiny partly} controlled by their parent
	\end{itemize}
\end{frame}

\begin{frame}
	\frametitle{Widgets in practice}
	\begin{itemize}
		\item These are only general guidelines
		\item Each widget has its own specificities
		\item Choose your widget using the screenshots (see \href{https://github.com/oliviercailloux/java-course/blob/master/SWT.adoc}{doc})
		\item Then check the javadoc or snippets
		\item Then check what JFace provides to help
	\end{itemize}
\end{frame}

\section{Practicalities}
\begin{frame}
	\frametitle{Practicalities}
	\begin{itemize}
		\item Maven: SWT not nicely mavenized (see \href{https://github.com/oliviercailloux/java-course/blob/master/SWT.adoc}{doc})
		\item UI Thread: the thread that created the display
		\item Access GUI widgets: only through the UI thread
		\item Event loop must be run from the UI thread
	\end{itemize}
\end{frame}

\appendix

\section{Licence}
\clearpage\pdfbookmark{License}{License}
\begin{frame}[plain]
	\frametitle{Licence}
	Cette présentation, et le code LaTeX associé, sont sous \href{http://opensource.org/licenses/MIT}{licence MIT}. Vous êtes libres de réutiliser des éléments de cette présentation, sous réserve de citer l’auteur.
	
	Le travail réutilisé est à attribuer à \href{http://www.lamsade.dauphine.fr/~ocailloux/}{Olivier Cailloux}, Université Paris-Dauphine.
\end{frame}
\addtocounter{framenumber}{-1}
\end{document}

\clearpage\pdfbookmark[2]{\refname}{\refname}
\begin{frame}[allowframebreaks]
	\frametitle{\refname}
% 	\bibliography{zotero}
\end{frame}

\section{Autres}
\begin{frame}
	\frametitle{}
	\begin{itemize}
		\item 
	\end{itemize}
\end{frame}

\end{document}
